
\sourcenote{Genre: core rule (tier 1). Setting as in the \rulename{Ax}
section; counterexamples read over the one-wire hypernet.}

\textbf{Claim (\texttt{sound\_andI}).}\enspace The payload supplies two
premise sequents $\Phi \vdash \Psi_1$ and $\Phi \vdash \Psi_2$ with
their sub-derivations and three tree identities: the two premises share
one context (\texttt{preMatch}), the conclusion's premise is that
context (\texttt{preEq}), and the conclusion's post is
$\Psi_1 \wedge \Psi_2$ with the left premise's conclusion as left
conjunct (\texttt{postEq}). If both premise sequents are semantically
valid, so is $\Phi \vdash \Psi_1 \wedge \Psi_2$.

\begin{bthm}[soundness of $\wedge$I]\label{thm:andi-sound}
Under the three payload identities, semantic validity of both premise
sequents implies semantic validity of the conclusion, over every
hypernet.
\end{bthm}
\begin{proof}
Let $\sigma$ be admissible with $\sigma \models \Phi$.

\emph{Step 1 (left component).} By \texttt{preEq}, $\sigma$ satisfies
the left premise's premise; left validity gives $\sigma \models \Psi_1$.

\emph{Step 2 (right component).} \texttt{preMatch} transports
$\sigma \models \Phi$ verbatim to the right premise's premise --- the
same tree, so the same satisfaction fact; right validity gives
$\sigma \models \Psi_2$.

\emph{Step 3 (pair).} The $\wedge$-clause of satisfaction is
meta-conjunction, so $\sigma \models \Psi_1 \wedge \Psi_2$, pinned by
\texttt{postEq}.
\end{proof}

\begin{bremark}[\texttt{preMatch} is essential]\label{rem:andi-share}
The rule may only pair conclusions drawn from the \emph{same}
assumptions --- context sharing again, as at \rulename{Cut}
\cite{girard1987}. Drop \texttt{preMatch} and take the left premise
$\top \vdash \top$ and the right premise $w \leq 0 \vdash w \leq 0$
(both \rulename{Ax} instances): the degenerate rule concludes
$\top \vdash \top \wedge (w \leq 0)$ --- false at $\sigma(w) = 1$,
because the right conjunct was earned under an assumption the
conclusion no longer carries. \texttt{preEq}'s removal admits the
context-swap family (the \rulename{Weaken} section); \texttt{postEq}'s
removal frees the conclusion outright.
\end{bremark}

\begin{bremark}[order and association are data]
\texttt{postEq} places the left premise's conclusion on the left of one
new $\wedge$-node. The conjunction is an ordered tree: nothing in the
core commutes or reassociates it, so $\Psi_1 \wedge \Psi_2$ and
$\Psi_2 \wedge \Psi_1$ are different conclusions, each derivable from
the same premises by choosing which sub-derivation rides which payload
slot.
\end{bremark}

\begin{bremark}[the meet adjunction, split across two rules]
\label{rem:andi-meet}
In the lattice reading, $\wedge$I is one direction of the universal
property of meet --- $x \leq y$ and $x \leq z$ iff
$x \leq y \sqcap z$ --- and $\wedge$E is the other. Together the two
rules say the $\wedge$-fragment of the calculus adds nothing beyond the
lattice structure of entailment; the $\vee$E section makes this precise
(and essential) through a countermodel over exactly that structure.
With $\wedge$E the pair is also computationally lossless: the two
projections undo pairing, and pairing the two projections rebuilds the
identical formula tree.
\end{bremark}

\begin{anchors}
Steps 1--3 & \texttt{sound\_andI}; payload \texttt{PAndIData} (fields \texttt{preMatch}, \texttt{preEq}, \texttt{postEq}); wired as \texttt{RuleData.andI}\\
$\wedge$-clause of satisfaction & \texttt{Formula.holds} (case \texttt{and})\\
lattice reading & \texttt{M3.andI\_kit}\\
\end{anchors}
